\section{The Seven Sacred Liberal Arts}

It is no coincidence that the Sacred Liberal Arts are Seven in number: this is the number of the mysteries, the steps towards union with God\footnote{\url{https://www.gornahoor.net/?p=10801}}, and the rays of Creation, as well as the visible planets, the notes in the musical scale, and the steps or levels of the Kabbalah. There are also seven deadly sins. We can collate some of them, as follows:

\begin{longtable}{*{5}{p{.15\textwidth}}}
	\caption{Correspondences between seven qualities}\\
	\toprule
	\textbf{Liberal Art} & \textbf{Augustine's Stage} & \textbf{Ray}\footnote{\url{https://www.theosophical.org/files/resources/books/SevenRays/SevenRAys.pdf}} & \textbf{Mystery} & \textbf{2 Peter}\footnote{\url{https://www.biblegateway.com/passage/?search=2+Peter+1\%3A5-7&version=KJV}} \textbf{Beatitudes} \\
	\midrule
	\endfirsthead
	
	\toprule
	\textbf{Liberal Art} & \textbf{Augustine's Stage} & \textbf{Ray} & \textbf{Mystery} & \textbf{2 Peter Beatitudes} \\
	\midrule
	\endhead
	
	\midrule
	\multicolumn{5}{r}{\textit{Continues in next page}}\\
	\endfoot
	
	\bottomrule
	\endlastfoot
	
	Grammar & Animation & Freedom and Action & Washing of Feet & Faith-Virtue \\
	Logic & Sentience & Unity and Philanthropy & Scourging & Virtue-Knowledge  \\
	Rhetoric & Rationality & Com\-pre\-hen\-sion and Philosophy & Crown of Thorns & Knowledge-Self Control  \\
	Arithmetic & Virtue & Harmony & Bearing the Cross & Self-Control-Steadfastness \\
	Geometry & Ataraxia & Truth and Scientist & Mystic Death & Steadfastness-Godliness \\
	Music & Intuition & Goodness and Religion & Entombment & Godliness-Brotherly Affection  \\
	Cosmology & Union with God & Beauty and Art & Resurrection & Brotherly Affection-Love \\
	
\end{longtable}


The liberal arts presuppose high intelligence, as the personal consciousness of the practitioner has to have already gotten as far as to see that there is no conflict, in the end, between Faith and Reason: since intellectus is destined for reunion with God just as surely as is purest faith. The union is effected in the consciousness of the just, or those made just. Those addicted to disputations or dilettantism will find the way of participation barred by the nature of the realities handled by the art. For those willing to re-engage them, the position and order of the seven liberal arts is teaching something about what its spiritual purpose is: it does this by the \emph{analogia entis\footnote{\url{http://www.millinerd.com/2006/12/whos-afraid-of-analogia-entis.html}}}, the law of correspondence, and the doctrine of signatures. The above table is a start.

For example, Mouravieff's Man \# 3 is a man of rhetoric (ideas); he engages the world intellectually, even when acting emotionally or instinctively, he uses the intellectual sectors with emotional or instinctual coloring heavily. Such a man would need to beware traps of rhetoric, and also study to engage more deeply in the truest and best use of creative rhetoric. Rationality, self-control, and the ``crown of thorns'' (the crucifixion of the intellect) would be things most necessary for this initiate to attend to in his self examination and training. He would attend to them knowing that they were his greatest strength, and also his greatest weakness, whereas weakness in the other disciplines would manifest completely differently: they would either be far more obvious, or far less -- A man \#3 would either be grossly deficient in the emotional arena, or (perhaps), far less tainted in that area (than in the intellect), and more artlessly or naturally pure, simply because it was not the seat of his ego, and precisely because it had been less ``developed''. Indeed, even in the exoteric literature, we find references to wisdom like this, as for example John Cardinal Newman recommending the benefits of liberal education\footnote{\url{http://www.newmanreader.org/works/idea/index.html}} for distracting the mind from idleness or temptation to fear, which is the besetting structural sin of a Man \#3.

Therefore although Geometry seems like a strange subject to study to develop one's spirituality, this is the ineluctable conclusion one is led to when delving into the interiority of man. For example, Thomas Hobbes came to Geometry\footnote{\url{https://www.uh.edu/engines/epi2372.htm}} rather late in life, and it shows in his thought, as illustrated in this amusing story:

\begin{quotex}
He was 40 years old before he looked on geometry; which happened accidentally. Being in a gentleman's library, Euclid's Elements lay open, and ``twas the 47 El. libri I'' Pythagoras' Theorem . He read the proposition . ``By God'', sayd he, ``this is impossible:'' So he reads the demonstration of it, which referred him back to such a proposition; which proposition he read. That referred him back to another, which he also read. Et sic deinceps, that at last he was demonstratively convinced of that trueth. This made him in love with geometry.

\flright{Quoted in \textsc{O L Dick}, \textit{Brief Lives} (Oxford 1960); A quotation about Thomas Hobbes by John Aubrey (1626-1697)}
\end{quotex}


Another observation that shows the non-arbitrary nature of the patterns in the Logos: the feet are the lowest part of the body, but also closely connected to spiritual development. Incidentally, this is the basis of ``foot-washing'' and also why people's feet get cold first when they are dying. There is also the old Bible verse: ``How beautiful upon the mountains are the feet of him that bringeth good tidings, that publisheth peace; that bringeth good tidings of good, that publisheth salvation; that saith unto Zion, Thy God reigneth''. (Isaiah 52:7) Mouravieff cites in Gnosis the practice developed by a group associated with Karel Weinfurter, in Man's Highest Development\footnote{\url{https://archive.org/details/manshighestpurpo032199mbp/page/n10}}, which details spiritual exercises designed to center development on the feet. The Middle Pillar Exercise\footnote{\url{https://www.golden-dawn.com/eu/UserFiles/en/file/pdf/Middle\_Pillar.pdf}} in Rosicrucian magic is a kind of ``grounding'', as well. Both exercises are a kind of magical ``Grammar'' (and Arithmetic). They are initiatory. And so can all the sacred liberal arts.

An entire program of spiritual and emotional and intellectual development is laid out in the Trivium and the Quadrivium, leading (finally) to union with God as its ultimate end. We can see that this seventh branch of the table is the highest plane of existence. The matching tables could make this very practical and specific. This seven fold pattern should be tried in a spirit of expectant humility: the numbers are not arbitrary.. He who maintains this out of sheer ignorance has not graduated the very first stage of instinctual life: the protozoa are more intelligent, in that they do better arithmetic with their instincts, and better grammar with their participation in Logos. The ignorant sceptic sorely needs the sacred liberal arts.

Recognizing the inner meaning of the liberal arts frees it from endless debates around the exact and minute particulars of the ``Canon'': it is less important what is precisely ``in'' (or ``out'') of the Canon, than in whether or not whatever is being studied shares a certain inner pattern and spirit. For example, \emph{Faust} is a Western classic. However, there are other obviously acceptable substitutes\footnote{\url{https://en.wikipedia.org/wiki/The_Tragedy_of_Man}} : someone can get an education from very few books, if they are ``right'' for their inner development. The development of depth in the pupil means a student of the \emph{Bhagavad Gita} and a disciple of the Gospel of John could hold a conversation and understand one another, because they share a deeper inner language. Canon details, provided they follow a pattern laid down by the Logos, are authentic and true. Esoteric recovery of the Liberal Arts will re-vitalize them, and save us from snobbery or dilettantism, or irrelevance.

Some will admit that the numerical pattern is interesting, but will find the details too speculative and abstruse. They are not cynical so much as possessing an over-developed worldliness. But their skepticism is even more difficult to address. Still, we can try to see if we can aid them, tossing a bone\footnote{\url{https://www.npr.org/templates/story/story.php?storyId=11803762}}:

\begin{quotex}
SIMON: I -- do you have any idea how seven became a lucky number?

Prof. DEVLIN: I have no idea, whatsoever, except that, you know, if we forget lucky numbers, there is a mathematical notion called a happy number. To get a happy number, you look at the number and you start squaring the digits and adding the answers. If you keep on doing that, eventually either you'll end up with the number one or you will end up getting a sequence of numbers that cycles. It begins 4, 16, 37. Seven is the smallest number that gives you one.
\end{quotex}


As you see, this is exactly the kind of stuff we find in Iamblichus' \textit{The Theology of Arithmetic}\footnote{\url{https://archive.org/stream/iamblichus-theologyarithmetic/iamblichus-theologyarithmetic_djvu.txt}}, which is dismissed today as risible. After leafing through the chapter on the Hebdomad (Aquinas commentated on Boethius' commentary\footnote{\url{https://www.amazon.com/Exposition-Hebdomads-Boethius-Aquinas-Translation/dp/0813209951}} on the Hebdomad), I did a little research, to see if I could find out what Iamblichus meant by calling 7 a virginal number. As it turns out, it is similar to seven being a ``happy'' number:

\begin{quotationx}
Evidence A\footnote{\url{https://books.google.com/books?id=udptayxY2E8C\&pg=PA27\&lpg=PA27\&dq=hebdomad+athena+virgin+iamblichus\&source=bl\&ots=2mH2dyjQgG\&sig=ACfU3U1fwbHMQmhzoZ67D8M2Hfwa4CJ_qw\&hl=en\&sa=X\&ved=2ahUKEwjQjZ60q-HhAhUNVa0KHTmLBI0Q6AEwC3oECAcQAQ\#v=onepage\&q=hebdomad\%20athena\%20virgin\%20iamblichus\&f=false}} (this source is derisive, but gives a valuable tidbit)

Other numbers arise naturally, but seven is ``left out'': 1-3 are generated by the point, the line, and the plane, and 4 is generated by the solid. 3 \& 4 then generate 5 in the Pythagorean triangle as hypotenuse, and 6 as the area, 8 is the first cube number, and 9 the first square. But Seven is nowhere to be found ``in mathematical nature'', except by intention.

Evidence B\footnote{\url{https://books.google.com/books?id=ifrHAAAAQBAJ\&pg=PA161\&lpg=PA161\&dq=virgin+seven+iamblichus\&source=bl\&ots=kef3A0I8kF\&sig=ACfU3U2Q5NAMgciQ-tERPOxlYWPxU5vmRA\&hl=en\&sa=X\&ved=2ahUKEwisq7e3ruHhAhUJvKwKHe0TAjUQ6AEwAHoECAkQAQ\#v=onepage\&q=virgin\%20seven\%20iamblichus\&f=false\%20target=}}

Agrippa states it is the only number that both 1) when multiplied by another number, produces none of the numbers in the decad (unlike 5, as in 5*2=10), and 2) is not the product of another multiplication of another number (only 1*7=7). So it has ``no father or mother''.
\end{quotationx}


Mathematics hasn't been affected by ``Modernity''. At least not in the qualities of the elemental numbers. The conclusion has to be that Seven is a repetition of One: It is Unity, Recreated and Manifested in fullness or perfection. Consult Genesis 1. Even the French Revolution had to go back to the days of the week.

Yet what other meaning could the numbers possess, if not to direct our attention to a particular spot, for deepening wisdom? For those lacking the arithmetical instincts of the amoeba or average honeybee because they deprive themselves, consider the alternative, which is that learning has no pattern, and no symbolism, but is purely utilitarian or even existential. Under this view, there is nothing requiring a leading out ``from'', and nowhere to quest ``towards''. Such positions all too obviously lead to nihilism, because if you destroy the symbolism of numbers, you are left with Zero. Others will find it somewhat fanciful, perhaps -- But why 7 and not 12? Or 21? What is the meaning of Life? You may as well answer 42\footnote{\url{https://www.independent.co.uk/life-style/history/42-the-answer-to-life-the-universe-and-everything-2205734.html}}! or else take Iamblichus seriously.

So certainly, each number ``means something''; they are something: thus, there are other numbers and patterns and meaning ``after'' 7, just as there are meanings for before 7. Since the intersection between Macrocosm (Triad) and Microcosm (Universe of ``4'') is 7, this has radical implications for education and spiritual development: our forefathers knew this in some way, whether instinctively, by experience, esoterically, or otherwise. Twelve (of course) results from the multiplication of 3 and 4, and is therefore just as significant as 7. And if we know how months differ from days, we have a clue to answer how. Perhaps the liturgical year could be brought into relation with the Liberal Arts? There are other things to say about the Numbers of course, all of them, and someone else could say more of 7. So the pattern can be enriched and expanded.

If God created the world by wounding himself and withdrawing to make the Void, then our imperfection is a task for both Him and our selves to overcome and finish the Creation, for God is at rest (He is perfect). He rests, having become man and overcome death, so that we might become God. God is, therefore, ``at work'' through us in our imperfection. He is perfect, and yet (on the human plane of activity), imperfect, since the revealing of the sons of God is not yet here:

\begin{quotex}
Praeterea, ostensum est supra quod essentia Dei est ipsum esse. Sed ipsum esse videtur esse imperfectissimum, cum sit communissimum, et recipiens omnium additiones. Ergo Deus est imperfectus. Objection 3: Further, as shown above (Question\footnote{\url{https://dhspriory.org/thomas/summa/FP/FP004.html}}3 , Article 4 ), God's essence is existence. But existence seems most imperfect, since it is most universal and receptive of all modification. Therefore God is imperfect.
\end{quotex}


Aquinas goes on to deny that God is imperfect (he answers this objection), yet there is an inner meaning to the old verse: And He said to me, ``My grace is sufficient for you, for My strength is made perfect in weakness (I Cor 12:8-10). What is this making perfect, if not God taking on imperfection, to suffer with those humble or pure enough to be of His mind, in order to further creative spiritual evolution?

Thus, the liberal arts are the strength of God for our weakness of mind, yet they are simultaneously nothing compared to the strength of God that they lead towards. For once completed, they are fulfilled: ``He who finds truth, beauty, and goodness, finds more'' (Dr. Michael Baumann). This is why Aquinas on his deathbed said\footnote{\url{http://scriptoriumdaily.com/thomas-aquinas-big-pile-of-straw/}}:

\begin{quotex}
I adjure you by the living almighty God, and by the faith you have in our order, and by charity that you strictly promise me you will never reveal in my lifetime what I tell you. Everything that I have written seems like straw to me compared to those things that I have seen and have been revealed to me.
\end{quotex}


Part of our task requires us to avail ourselves of a rational and sound mind, and to think God's thoughts after Him. In order to transcend humanity, we have to develop our humanity up to its natural limit: the common saying in parlance is, Put Wings to your Prayers. Trust in God, tie up your camel.

Sharpening the mind in, through, and on the sacred liberal arts, allows the fastest progress in developing the intellect up to its natural limit.

Hopefully, more work will be done in recovering the basis of the Seven Sacred Arts. The orders of angels comes to mind, as does Bonaventura's Journey of the Mind into God\footnote{\url{http://www.ntslibrary.com/PDF\%20Books/Bonaventure\%20Journey\%20of\%20the\%20Mind\%20Into\%20God.pdf}}: for example, if you add philosophy and theology to the seven liberal arts, you have 9, which is the order of the angels. A young scholar aspiring to be a sage could do worse than undertake the revitalization of the Sacred Seven Liberal Arts. Interest in it is already increasing, as evinced by this recent article\footnote{\url{http://www.tomworrel.com/uploads/3/4/7/9/34794642/ahiman_worrel_article_page_7-37.pdf}}. The table provided at the start, along with the esoteric background, may be of very practical use. Our forefathers, even a hundred years ago, had not entirely forgotten that the edifice of civilization is built, not on Mammon, but on something much deeper. They even carved it into stone\footnote{\url{http://www.studiesincomparativereligion.com/uploads/ArticlePDFs/403.pdf}}, so we could not forget.

Hopefully, we have scratched the surface enough in the above to demonstrate that there is a deep organic inner Logic behind the tradition and lay out of the Seven Sacred Liberal Arts, enough to get you interested, and started. It's just a beginning.


\flrightit{Posted on 2019-05-05 by Logres}

\begin{center}* * *\end{center}

\begin{footnotesize}
\begin{sffamily}
\texttt{Zach on 2019-05-09 at 21:15 said:}

Thanks. This came on a significant day for me.

Although I am well past being a ``young scholar.'' The Hobbes anecdote is hopeful.

\hfill

\texttt{Logres on 2019-05-12 at 00:20 said:}

I like the quote that Cologero put up recently, that the young learn from us what we've learned by heart (I can't remember the name of the Frenchman quoted). Falling in love with it, or all over again, helps us learn it by heart. Maybe you can help pass it on. Thanks for reading.

\end{sffamily}
\end{footnotesize}

